% In this file you should put the actual content of the blueprint.
% It will be used both by the web and the print version.
% It should *not* include the \begin{document}
%
% If you want to split the blueprint content into several files then
% the current file can be a simple sequence of \input. Otherwise It
% can start with a \section or \chapter for instance.
\chapter{Definitions}

\section{Introduction}

De Bruijn indices provide a way to represent lambda calculus terms without naming bound variables.
Instead of using variable names, each variable is represented by a natural number that indicates how many binders \( \lambda \) separate it from its corresponding abstraction.

\begin{definition}
  \label{def:Lambda}
  \lean{Lambda}
  \leanok
  Formally, the syntax of lambda terms using de Bruijn indices is defined as:
  \[
    \Lambda := \mathbb{N} \mid (\Lambda\ \Lambda) \mid (\lambda\ \Lambda).
  \]
  Here, variables are natural numbers, and each number refers to its binder according to its distance.
\end{definition}

\begin{example}
  For example, the lambda term \( \lambda x. \lambda y. x y \) is represented using de Bruijn indices as \( \lambda \lambda 1 0 \).
  In this representation, the variable \( x \) is represented by the index \( 1 \) because it is one binder away from its abstraction, while \( y \) is represented by \( 0 \) since it is directly bound by the nearest \( \lambda \).
\end{example}

The shift operation is defined to adjust the indices of variables in a lambda term.
Given a term \( N \) and naturals number \( d \) and \( c \), the shift operation \( \uparrow_c^d N \)
increases the indices of all free variables in \( N \) whose values are greater than or equal to the
context \( c \) by \( d \).

\begin{definition}[shifting]
  \label{def:shift}
  \lean{Lambda.shift}
  \uses{def:Lambda}
  \leanok
  Formally, the shift operation is defined recursively as follows:
  \[
    \begin{aligned}
      \uparrow_c^d n &=
      \begin{cases}
        n & \text{if } n < c, \\
        n + d & \text{otherwise},
      \end{cases} \\[1em]
      \uparrow_c^d(\lambda N) &= \lambda (\uparrow_{c + 1}^d N), \\[0.5em]
      \uparrow_c^d(N_1\ N_2) &= (\uparrow_c^d N_1)\ (\uparrow_c^d N_2).
    \end{aligned}
  \]
\end{definition}

Here, \(c\) is called the \textbf{context}, and \(d\) the \textbf{displacement}.
Similarly, we define the \textbf{unshifting} operation, which behaves the same
way except that it decreases the index by $d$ instead of increasing it:

\begin{definition}[unshifting]
  \label{def:unshift}
  \lean{Lambda.unshift}
  \uses{def:Lambda}
  \leanok
  \[
    \begin{aligned}
      \downarrow_c^d n &=
      \begin{cases}
        n & \text{if } n < c, \\
        n - d & \text{otherwise},
      \end{cases} \\[1em]
      \downarrow_c^d (\lambda N) &= \lambda (\downarrow_{c + 1}^d N), \\[0.5em]
      \downarrow_c^d (N_1\ N_2) &= (\downarrow_c^d N_1)\ (\downarrow_c^d N_2).
    \end{aligned}
  \]
\end{definition}

\section{Substitution}

\begin{definition}[Substitution]
  \label{def:substitution}
  \lean{Lambda.subst}
  \uses{def:Lambda,def:shift}
  \leanok
  The substitution operation replaces all occurrences of a variable in a term with another term.
  Its formal definition is as follows:
  \[
    \begin{aligned}
      n[m := M]&=
      \begin{cases}
        M & \text{if } n = m, \\
        n & \text{otherwise},
      \end{cases} \\[1em]
      (\lambda N)[m := M] &= \lambda N[m + 1 := \uparrow_0^1 M], \\[0.5em]
      (N_1\ N_2)[m := M] &= (N_1 [m := M]) (N_2 [m := M]).
    \end{aligned}
  \]
\end{definition}

In the case of an abstraction, we move one level deeper into the binding structure, so the context increases by one. That's why we use \(m + 1\).

Additionally, we apply shifting to the term \(M\) being substituted. This is necessary because \( M \) is inserted into a context with one extra binder. Without shifting, the free variables of \(N\) could become incorrectly bound, changing the meaning of the expression.
